\section{What Was Not Tested}
\label{sec:untested}

This is stated up front and in one place, rather than left to be inferred from the
absence of a citation in Section~\ref{sec:traceability} or from the environment of
Section~\ref{sec:environment}.  None of the items below is hidden elsewhere in this
document; they are gathered here so a reader does not have to hunt for them.

\subsection{Platforms and browsers never run}

\begin{itemize}
  \item \textbf{Operating systems other than the one in Section~\ref{sec:environment}.}
        macOS is the only supported platform (\fq{DST-4} is macOS-specific by design),
        but even within that, only one build and one architecture were tested.  Intel
        Macs, and Apple-silicon Macs other than the model listed, have not run any part
        of this program.  On Linux, the Go tests and the Chrome end-to-end suite run in
        continuous integration (on GitHub's hosted runners), but \texttt{fsb} has never
        been used there by hand, and Quick Look previews do not exist there.
  \item \textbf{Firefox.}  Not installed on the test machine, not driven by any
        end-to-end test, and not opened by hand.  The interface uses no
        Chrome-only or Safari-only API knowingly, but this is design intent, not a
        verified fact for Firefox.
  \item \textbf{Any browser's mobile mode, and any touch or narrow-window
        interaction.}  The narrow-window column behavior (\fq{LST-11}) was written to
        a requirement, not exercised at a phone-sized viewport.
  \item \textbf{A second machine, a second user account, or a network-mounted
        volume.}  Every test ran as one user on one machine.  Cross-user permission
        boundaries (distinct from the deny rules) are provided by the operating system
        and were not separately exercised.
\end{itemize}

\subsection{Scenarios inside the test suite that did not run this time}

\begin{itemize}
  \item \textbf{The 100{,}000-file scale test} (\fq{LST-2}) is opt-in
        (\code{FSB\_E2E\_BIG=1}).  It runs in continuous integration, on Linux, where the
        first run (2026-10-02) drew the first rows of 100{,}000 entries after 118~ms with
        35 rows in the page.  On macOS it was run by hand once, on 2026-09-20, which is
        what the functional specification's limits section reports.
  \item \textbf{\code{TestRootContainmentIsByIdentityNotSpelling}} (\fq{ROOT-1}) is
        skipped unless the environment variable \code{FSB\_CASE\_SENSITIVE\_VOLUME}
        names a folder on a case-sensitive volume.  It was run once, manually, against a
        disk image created for that purpose (\code{hdiutil create -fs "Case-sensitive
        APFS"}), and is skipped in every other run, including in CI.
  \item \textbf{Real cloud-only (``dataless'') files and folders.}  \fq{PRV-10},
        \fq{LST-13} and \fq{SEC-13} are tested against a synthetic dataless flag set on
        an ordinary local file or folder (\code{TestIsDataless}, \code{TestRefuseDataless}
        and their callers), because the real flag is set by the File Provider framework
        itself and cannot be created to order: setting it directly with \code{chflags}
        was tried and confirmed not to work (it reports success but the flag does not
        persist), and a real iCloud Drive or OneDrive placeholder cannot be created in a
        test fixture or a CI runner.  Findings 19 and 20 (Section~\ref{sec:history}) are
        exactly this gap: a real File Provider folder was the only thing that ever
        exercised the code path that was missing.  When fixing them, reaching the actual
        reported folder to confirm the fix live was also tried and abandoned: the
        process running the fix could not read \code{\textasciitilde/Library/CloudStorage}
        at all (macOS refused it outright), a permission boundary distinct from the one
        the user's own session had already been granted, so even this attempt at direct
        confirmation was itself an instance of the same limitation.
  \item \textbf{A real mouse drag of a column header in Safari.}  The Safari driver
        cannot dispatch a native drag-and-drop sequence, so the browser test
        (\fq{LST-11}) falls back to dispatching the underlying drag events directly
        (logged by the suite as \emph{``safari did not complete a native drop \ldots;
        falling back to synthetic drag events''}); this has not been confirmed with an
        actual pointer drag in Safari.
  \item \textbf{Real cloud-only iWork and Office documents} (\fq{PRV-12}).  The
        refusal is the same check as every other preview's and is tested the same
        synthetic way; no real cloud-only document has been previewed.
  \item \textbf{Encrypted PDFs} (\fq{PRV-8}).  The test moves the focus into the PDF
        frame itself; a real encrypted PDF, whose viewer takes the focus, was tried only
        by hand (by the user, 2026-10-02).
  \item \textbf{The clipboard} (\fq{LST-9}, Alt+C).  The end-to-end driver
        cannot read the system clipboard back in every environment (logged as
        \emph{``clipboard contents not read back with this driver''}), so the copy
        action is exercised but the copied text is not always verified.
\end{itemize}

\subsection{Process and infrastructure}

\begin{itemize}
  \item \label{sec:ci-never-run}\textbf{Continuous integration is new.}  It first ran on
        2026-10-02, when the repository was given a private remote; that run found three
        problems (finding 28 of Section~\ref{sec:history}).  It runs the Go tests on
        macOS and Linux, the front-end tests, fuzzing, \code{govulncheck}, \code{npm
        audit}, and the Chrome suite on Linux, on every push.
  \item \textbf{Safari's end-to-end suite does not run in CI.}
        \code{ci.yml} sets \code{FSB\_E2E\_SKIP\_SAFARI=1} for the hosted job, because a
        headless CI runner cannot drive Safari.  The Safari suite has been run only on
        this machine, by hand, each time a feature changed.
  \item \textbf{No independent review has been performed by anyone other than the
        author and an assistant.}  The two ``independent reviews'' recorded in
        Section~\ref{sec:history} and \code{specs/SECURITY.md} were done by an assistant
        with access to the source but none to the author's reasoning, not by a third
        party, and not by a professional security firm.  No penetration test by an
        independent party was done; the owner decided against one.
  \item \textbf{Installing on another machine.}  The release is source only
        (\fq{DST-2}), so there is no signing, notarization or Gatekeeper behavior to test,
        but \code{go install} from the public repository was tried only once, on the
        machine of Section~\ref{sec:environment}, right after release.
  \item \textbf{Vulnerability scanning.}  \code{govulncheck} and \code{npm audit} run in
        continuous integration on every push, and GitHub's Dependabot alerts flag a
        dependency when a new advisory names it.  No automated dependency updates are
        proposed, and nobody is committed to acting on an alert.
\end{itemize}

\subsection{What this does not mean}

None of the above is a reason to distrust the parts that \emph{were} tested: the
guard, the rule matcher, the HTTP middleware and the endpoints are exercised by the
heaviest layer of testing this project has (Section~\ref{sec:verification}), including
a differential test that asks the operating system itself, not \texttt{fsb}, which
paths are the same file.  The gaps above are concentrated in exactly the places that
testing is normally weakest anywhere: process (CI is new), rare environments
(a case-sensitive volume, real cloud placeholders, another OS), and interactions a
test harness cannot fully simulate (a native drag, a system clipboard).  They are
listed here so that a reader, in particular someone doing a pre-release review, knows
where to look first.
